% Paragraph 1 (both storylines): classical result and the central problem.
% Unchanged from Introduction_final.tex.
Cooperation among unrelated individuals is sustained in part by reputation\cite{milinski2016,nowak2005}. By helping others, individuals improve their standing and become more likely to receive help from third parties, a mechanism known as indirect reciprocity\cite{nowak1998,leimar2001,wedekind2000,milinski2002}. Theory has identified social norms that make reputation-based cooperation evolutionarily stable\cite{panchanathan2003,ohtsuki2004}, most notably the leading-eight norms, which sustain cooperation when reputations are public, even in the presence of occasional errors\cite{ohtsuki2004,ohtsuki2006}. These results rest on agreement about who is good and who is bad. When reputations are private and assessments are noisy, individuals hold different views of the same person\cite{takahashi2006,uchida2010}, and most leading-eight norms no longer maintain stable cooperation\cite{hilbe2018,okada2018,perret2021,fujimoto2023}. A central problem for indirect reciprocity is therefore how reputations can remain both accurate and shared under private assessment. They must distinguish those who cooperate from those who do not, and different observers must largely agree about who is who\cite{fujimoto2022,murase2024synchronization}.
